\documentclass[11pt]{article}
\usepackage[a4paper,margin=2cm]{geometry}
\usepackage{amsmath,amssymb,amsthm}
\usepackage{listings}
\usepackage{xcolor}
\usepackage{hyperref}
\definecolor{critical}{rgb}{0.8, 0.1, 0.1}

\title{\textcolor{critical}{\textbf{[CRITICAL]}} Plan:
  HymeKo-Nagare $\to$ HSMM Compiler\\
  \emph{Future-session work; FPGA prerequisite}}
\author{}
\date{Drafted: 2026-05-11. Status: \textbf{frozen}, awaiting FPGA HSMM kernel.}

\begin{document}
\maketitle

\section*{\textcolor{critical}{Status banner}}

This is a \textbf{CRITICAL} plan flagged for a \textbf{future session}.
It is \textbf{not} actionable in the current session because:

\begin{enumerate}
  \item The \textbf{HSMM FPGA implementation is a prerequisite}.
    HSMM (Hypergraph Storage Modification Machine, Hajdu \& Csapó 2026
    Zenodo manuscript) currently has a Zynq UltraScale+ kernel
    referenced in \texttt{docs/plans/plans\_20260429/hymeko\_monitor\_plan.md}
    but the bitstream / HDL artefacts are not in this repo
    (verified via \texttt{grep}). The compiler target must exist
    before the source-language compiler is meaningful.
  \item The Nagare ops landed today are CPU-only; GPU backend
    (Stage 3 of \texttt{docs/plans/2026-05-11-hymeko-nagare-flow/plan.tex})
    has not yet been built. Going straight to FPGA would skip a
    natural staging point.
  \item Empirical validation requires a Nagare $\to$ HSMM
    end-to-end pipeline that we can compare against PyTorch-on-CPU
    on real signed-graph workloads --- requires the FPGA target
    booted and the Nagare ops mature.
\end{enumerate}

\textbf{Re-open this plan when:} the HSMM FPGA bitstream is in repo
(\texttt{hymeko\_monitor/hsmm\_bridge.rs} mentions it; needs to
materialise), OR the Hajdu-Csapó companion paper publishes a public
HSMM reference implementation, whichever comes first.

\section{Goal}

Build a compiler $\phi : \text{Nagare-program} \to \text{HSMM-program}$
such that every Nagare training-loop function (forward + closed-form
backward + Adam step + loss) compiles to a fixed HSMM-program
(sequence of NEWE / ATT / arithmetic primitives) executable on the
Zynq UltraScale+ HSMM kernel.

This is the systems counterpart to the algorithmic paper
\texttt{docs/plans/2026-05-11-hymeko-nagare-flow/plan.tex}: it
demonstrates that Nagare is the natural ML programming language for
HSMM (Schönhage-class storage-modification machine, hypergraph
variant), via a structural embedding rather than a trace-based
lowering.

\section{Why this matters}

\textbf{For PyTorch / JAX / Burn:} reverse-mode autograd needs to
traverse the recorded computation graph during backward. Compiling
that to HSMM would require lowering the graph traversal itself to
HSMM primitives --- effectively building an HSMM-resident autograd
interpreter. That is feasible but bypasses the whole point of
HSMM (compiled, fixed-schedule execution).

\textbf{For Nagare:} the closed-form Clifford backward (Theorem 1 of
\texttt{math.tex} in the Nagare plan dir) is already expressed as a
fixed sequence of READ / FMA / ATT-LOCAL operations over the
hyperedge stream. There is no graph traversal at backward time. The
compilation to HSMM is a \emph{structural} embedding (operator
$\to$ template), not a \emph{trace-based} one. This is the property
that makes the Nagare $\to$ HSMM pairing structurally clean.

\section{Affected files (when reopened)}

\begin{lstlisting}
hymeko_nagare/
├── src/
│   ├── compiler/
│   │   ├── mod.rs             pub use of the compiler
│   │   ├── ir.rs              Nagare op IR (op-graph in flat Vec form)
│   │   ├── hsmm_target.rs     HSMM program emit (NEWE / ATT / FMA)
│   │   ├── schedule.rs        operator-level scheduling for FPGA pipeline
│   │   └── verify.rs          structural-equivalence checker (Nagare op
│   │                          vs emitted HSMM program)
│   └── ...                    (existing ops crate)
hymeko_monitor/
└── src/
    └── hsmm_bridge.rs         existing scaffold; will host the runtime
                                bridge between Nagare-compiled HSMM
                                programs and live FPGA events.
docs/plans/CRITICAL-nagare-to-hsmm-compiler/
├── plan.tex (this doc)
├── plan.pdf
├── plan.tikz                  Nagare-IR -> HSMM-program structural diagram
├── plan.mmd                   compiler stages flowchart
└── README.md                  prerequisites checklist + reopen criteria
\end{lstlisting}

\section{Prerequisites checklist}

Before reopening this plan, ALL must be true:

\begin{enumerate}
  \item HSMM Zynq bitstream OR public reference implementation
    available in repo, with documented NEWE / ATT primitive
    surface (event format, memory layout, dispatch protocol).
  \item Nagare Stage 3 (GPU backend via wgpu OR cudarc) shipped,
    so we have a CPU $\to$ GPU step \emph{before} CPU $\to$ FPGA;
    that gives us 2 architectural-portability validation points
    before the FPGA attempt.
  \item Catmull-Rom spline forward + closed-form backward ported
    to pure Rust (Nagare Stage 4), so the spline-aggregator path
    is also compilable to HSMM. Without this, only the FIR path
    works on HSMM, and the comparison vs the full Slashdot
    SOTA recipe is not meaningful.
  \item Slashdot 5-seed paired AUC reproduced on Nagare-CPU
    pipeline at parity ($\le 0.005$ AUC drift) vs the existing
    PyTorch + Triton kernel pipeline. This validates that Nagare's
    pure-Rust ops have no accuracy regression before attempting
    FPGA deployment.
\end{enumerate}

\section{Compilation stages (when reopened)}

\begin{enumerate}
  \item \textbf{Parse Nagare op chain.} Take a Rust function of
    type \texttt{fn(\&CyclePool, \&Features) -> \&Logits} composed
    of registered Nagare ops (no autograd traces, no closures with
    captured state). Produce a flat Nagare-IR.
  \item \textbf{Verify closure under op registry.} Every node in
    the IR must be a registered Nagare operator with declared
    forward + closed-form backward kernels. Reject if any op is
    autograd-traced (which would require an interpreter on FPGA).
  \item \textbf{Lower to HSMM-IR.} For each Nagare op, instantiate
    the corresponding HSMM-template (per the structural embedding
    $\phi$ defined in \texttt{math.tex} Section 7). Templates are
    parameterised by (cycle batch size, feature dim, filter length,
    optional Highway/attention sub-templates if Stage 4 has landed).
  \item \textbf{Schedule for FPGA pipeline.} Assign operators to
    FPGA pipeline stages; pipeline-balance via fanout / fanin
    blockers; check resource budget (DSP slices, BRAM tiles,
    LUT-FF pairs) against the Zynq UltraScale+ chip-specific
    budget.
  \item \textbf{Emit Verilog / SystemVerilog} via the existing
    \texttt{hymeko\_monitor} LIR $\to$ SystemVerilog path (see
    plan \texttt{docs/plans/plans\_20260429/hymeko\_monitor\_plan.md}).
    The Nagare compiler reuses this emission step --- it just
    emits a different LIR program shape.
  \item \textbf{Bitstream synthesis + deploy.} Standard Xilinx
    Vivado pipeline. Out of scope for the Nagare compiler crate,
    in scope for the integration tests.
  \item \textbf{Empirical validation.} Run Slashdot or Epinions
    forward + backward pass on the FPGA-resident compiled Nagare
    program; report throughput (cycles/s, edges/s) and compare
    against PyTorch CPU baseline + Nagare CPU baseline. Expected
    range: $10$--$100\times$ over PyTorch CPU at iso-FLOPS, with
    much lower energy/op (FPGA-typical $10$--$50\times$ better
    than CPU at signal-processing-class workloads, which the
    FIR + scatter primitives are).
\end{enumerate}

\section{Test strategy}

\textbf{Compiler unit tests} (pure-CPU, no FPGA):
\begin{itemize}
  \item Each Nagare-op $\to$ HSMM-template mapping has a unit
    test that emits the HSMM program and verifies it against a
    hand-written reference HSMM program for that op.
  \item Round-trip: Nagare program $\to$ HSMM-IR $\to$ executed
    in a Rust HSMM-IR interpreter $\to$ result matches Nagare
    CPU forward to $\le 10^{-6}$.
  \item Backward parity: emitted HSMM program for backward
    matches Nagare CPU closed-form backward to $\le 10^{-6}$.
\end{itemize}

\textbf{FPGA integration tests} (requires bitstream on Zynq):
\begin{itemize}
  \item Slashdot k=3 cycle pool, forward only, output parity
    $\le 10^{-5}$ vs Nagare CPU. (Float-precision drift from FPGA
    fixed-point arithmetic, if used, is the limiting factor.)
  \item Full training loop (forward + bwd + Adam step) for 10
    epochs, AUC parity $\le 0.01$ vs Nagare CPU.
  \item Throughput measurement: edges processed per second; target
    $\ge 10\times$ PyTorch CPU at iso-FLOPS.
\end{itemize}

\section{Performance budget}

\begin{itemize}
  \item Compile time: $\le 30$ s for a typical Nagare training
    loop (parsing + lowering + verilog emission). FPGA synthesis
    via Vivado is hours; that's not the compiler's wall.
  \item FPGA latency: $\le 1$ ms per forward pass on Bitcoin OTC
    scale (30K cycles, k=3, d=32); proportionally for larger.
  \item FPGA energy: $\le 1$ W per training step (compared to
    $50$--$100$ W for the equivalent on a 7nm CPU). The energy-cost
    argument is the headline for the deployment / edge-inference
    pitch.
\end{itemize}

\section{Rollback path}

If FPGA deployment proves infeasible for any reason (Vivado bug,
resource budget overrun, fixed-point precision unacceptable for
AUC parity), Nagare's CPU + wgpu/cudarc backends remain fully
functional. The compiler crate is a separate module; deleting it
does not regress the Nagare CPU+GPU path. Worst case: the FPGA
compiler is shelved indefinitely, Nagare ships as a CPU+GPU
framework, and the Hajdu-Csapó HSMM paper stands on its own
(without the Nagare integration as proof-of-feasibility).

\section{Risk anticipation}

\begin{itemize}
  \item \textbf{Fixed-point precision.} FPGAs typically use
    fixed-point arithmetic; f32 emulation eats DSP slices.
    AUC parity at fixed-point Q16.16 vs f32 needs early validation
    (Stage 6 in the compilation pipeline).
  \item \textbf{HSMM primitive surface drift.} The NEWE/ATT
    primitive set in the Hajdu-Csapó manuscript may evolve before
    publication. The compiler should target an explicitly versioned
    HSMM-IR (\texttt{hsmm\_v1.0}) so the source-language remains
    stable when the target evolves.
  \item \textbf{Scope creep into attention.} HSiKAN's
    quaternion-attention edge-cr Highway path has a more complex
    closed-form backward; whether to support attention in the
    Nagare compiler depends on whether attention proves necessary
    for SOTA-on-FPGA. Default: ship the FIR-only path first.
  \item \textbf{HSMM Zynq performance ceiling.} Zynq UltraScale+
    has $\sim 2000$ DSP slices; large cycle pools (Epinions $530$K
    cycles) may exceed throughput at d=32. Mitigated by streaming
    the cycle pool in chunks of $10$--$50$K cycles (the existing
    \texttt{HSIKAN\_CYCLE\_BATCH=2000} pattern), which fits the
    DSP budget comfortably.
\end{itemize}

\section{Venue strategy}

This work, when shipped, completes a three-paper arc:

\begin{enumerate}
  \item \textbf{Hajdu \& Csapó 2026 (Zenodo, in review):} HSMM as
    abstract computational model. Pure theory.
  \item \textbf{Nagare paper (planned 2026, MLSys 2027):} Cycle-additive
    operators with closed-form Clifford backward; structural
    correspondence to HSMM via embedding $\phi$. Theory + CPU
    implementation.
  \item \textbf{This plan (Nagare-to-HSMM compiler, planned 2027):}
    Empirical realisation: compiled Nagare programs running on
    Zynq UltraScale+ HSMM kernel, with throughput + energy
    measurements vs PyTorch CPU + Nagare GPU baselines. Target
    venues: FPGA 2028 / FCCM 2028 / DAC 2028 / FCCM 2028; or
    ASPLOS 2028 if the framing leans architecture rather than
    EDA.
\end{enumerate}

The three papers together establish HSMM as a viable parallel-ML
substrate with a working software stack and empirical advantage,
a much stronger claim than HSMM theory alone.

\section{Citation arc}

\begin{itemize}
  \item Schönhage 1980 (original Storage Modification Machine)
  \item Knuth-Tarjan pointer machines (1979)
  \item Kolmogorov machines (1953)
  \item Hajdu \& Csapó 2026 (HSMM = hypergraph SMM, Zenodo)
  \item HymeKo-Nagare 2026 (\texttt{docs/plans/2026-05-11-hymeko-nagare-flow/})
  \item This plan (Nagare $\to$ HSMM compiler, 2027 work)
\end{itemize}

\end{document}
